% ---- pg-473 （书页 461） ----

费希尔的线性判别函数是在假定两个总体有相同协方差矩阵的条件下得到的. 因此毫不奇怪, 费希尔方法与最小期望错分代价法则这种特殊情况相对应. 分类法则(11-18)中的第一项 $\hat{y}=(\bar{\mathbf{x}}_1-\bar{\mathbf{x}}_2)'\mathbf{S}_{\mathrm{p}}^{-1}\mathbf{x}$ 就是费希尔的线性函数, 它使“组间变差”相对“组内变差”而言达到最大[见\hyperref[eq:11-23]{式(11-23)}]. 整个表达式
\phantomsection\label{eq:11-26}
\begin{equation}
\begin{aligned}
\hat{w}&=(\bar{\mathbf{x}}_1-\bar{\mathbf{x}}_2)'\mathbf{S}_{\mathrm{p}}^{-1}\mathbf{x}-\frac{1}{2}(\bar{\mathbf{x}}_1-\bar{\mathbf{x}}_2)'\mathbf{S}_{\mathrm{p}}^{-1}(\bar{\mathbf{x}}_1+\bar{\mathbf{x}}_2)\\
&=(\bar{\mathbf{x}}_1-\bar{\mathbf{x}}_2)'\mathbf{S}_{\mathrm{p}}^{-1}\left[\mathbf{x}-\frac{1}{2}(\bar{\mathbf{x}}_1+\bar{\mathbf{x}}_2)\right]
\end{aligned}
\tag{11-26}
\end{equation}
常被称为安德森(Anderson)分类函数(或称安德森统计量). 若 $[(c(1\mid2)/c(2\mid1))(p_2/p_1)]=1$, 则 $\ln[(c(1\mid2)/c(2\mid1))(p_2/p_1)]=0$, 法则(11-18)与基于费希尔线性判别函数的法则(11-26)相似. 这样, 费希尔分类法则等价于在相同先验概率和相同错分代价下的最小 ECM 法则.

\mvsub{分类是好的思想吗}

简而言之, 对两个总体, 通过考虑多元观测值的线性组合所得到的最大相对分离度等于距离 $D^2$. 这样就很方便, 因为在某些情况下, $D^2$ 可用于检验两个总体均值 $\boldsymbol{\mu}_1$ 和 $\boldsymbol{\mu}_2$ 是否有显著差异. 因此, 均值向量间有无差异的检验可看成是对所能达到的总体分离度的显著性检验.

假定总体 $\pi_1$ 和 $\pi_2$ 服从有相同协方差矩阵 $\boldsymbol{\Sigma}$ 的正态分布, 则如同在 6.3 节中那样, 对假设 $H_0:\boldsymbol{\mu}_1=\boldsymbol{\mu}_2$ 和 $H_1:\boldsymbol{\mu}_1\neq\boldsymbol{\mu}_2$ 的检验可用统计量
\[
\left(\frac{n_1+n_2-p-1}{(n_1+n_2-2)p}\right)\left(\frac{n_1n_2}{n_1+n_2}\right)D^2
\]
来完成, 这个统计量服从 $F$ 分布, 自由度为 $\nu_1=p$, $\nu_2=n_1+n_2-p-1$. 如果 $H_0$ 被拒绝, 我们可以认为 $\pi_1$ 与 $\pi_2$ 之间的分离度是显著的.

\noindent\mvbold{评注}\quad 显著分离并不一定意味良好的分类. 正如我们在 11.4 节中将会见到的, 可以对一个分类方法的有效性作出评估而不必依赖任何分离性检验. 另一方面, 如果分离不显著, 则寻找有用的分类法则将被证明是徒劳的.

\mvsub{当 \texorpdfstring{$\boldsymbol{\Sigma}_1\neq\boldsymbol{\Sigma}_2$}{Sigma1 != Sigma2} 时正态总体的分类}

如同所预料的那样, 当总体协方差矩阵不同时, 分类法则就更复杂.

考虑\hyperref[eq:11-10]{式(11-10)}中的多元正态密度, 并用 $\boldsymbol{\Sigma}_i$($i=1,2$)代替其中的 $\boldsymbol{\Sigma}$. 这样, 对两个总体来说, 不仅两个均值向量, 就是两个协方差矩阵也都不相同. 正如我们所见到, 最小 ECM 和最小 TPM 的区域依赖于密度比 $f_1(\mathbf{x})/f_2(\mathbf{x})$, 或等价地依赖于密度比的自然对数 $\ln[f_1(\mathbf{x})/f_2(\mathbf{x})]=\ln[f_1(\mathbf{x})]-\ln[f_2(\mathbf{x})]$. 当多元正态密度有不同的协方差结构时, 密度比中包含 $|\boldsymbol{\Sigma}_i|^{1/2}$ 的项就不会像 $\boldsymbol{\Sigma}_1=\boldsymbol{\Sigma}_2$ 时那样可以消掉. 此外, $f_1(\mathbf{x})$ 和 $f_2(\mathbf{x})$ 指数中的二次型也不能合并成如\hyperref[eq:11-13]{式(11-13)}中那样相当简单的结果.

将具有不同协方差矩阵的正态密度代入\hyperref[eq:11-6]{式(11-6)}, 取自然对数再化简(见练习 11.15), 得到分类域:
\phantomsection\label{eq:11-27}
\begin{equation}
\begin{aligned}
R_1:\ &-\frac{1}{2}\mathbf{x}'(\boldsymbol{\Sigma}_1^{-1}-\boldsymbol{\Sigma}_2^{-1})\mathbf{x}+(\boldsymbol{\mu}_1'\boldsymbol{\Sigma}_1^{-1}-\boldsymbol{\mu}_2'\boldsymbol{\Sigma}_2^{-1})\mathbf{x}-k\geqslant\ln\left[\left(\frac{c(1\mid2)}{c(2\mid1)}\right)\left(\frac{p_2}{p_1}\right)\right]\\
R_2:\ &-\frac{1}{2}\mathbf{x}'(\boldsymbol{\Sigma}_1^{-1}-\boldsymbol{\Sigma}_2^{-1})\mathbf{x}+(\boldsymbol{\mu}_1'\boldsymbol{\Sigma}_1^{-1}-\boldsymbol{\mu}_2'\boldsymbol{\Sigma}_2^{-1})\mathbf{x}-k<\ln\left[\left(\frac{c(1\mid2)}{c(2\mid1)}\right)\left(\frac{p_2}{p_1}\right)\right]
\end{aligned}
\tag{11-27}
\end{equation}
其中
